\section{第\ref{sec:neurontypes}章　ニューロンの種類}

この章の数値は、ヒトの脳で細胞の直接計数があるところではそれに基づく。「ニューロン1個に神経伝達物質1種類」という規則は、細胞アトラスと、例外も見つけた実験に基づく。\nt{DOPA}の役割はスパイクの記録に基づき、その他のブロードキャストチャネルの役割は仮説に基づく。

{\footnotesize
\setlength{\tabcolsep}{3pt}\renewcommand{\arraystretch}{1.15}
\begin{longtable}{>{\raggedright}p{0.5cm}>{\raggedright}p{4.0cm}>{\raggedright}p{5.6cm}>{\raggedright}p{2.3cm}>{\raggedright\arraybackslash}p{1.7cm}}
\toprule
No. & 主張 & 実験と測定されたもの & 出典 & 状態 \\
\midrule\endfirsthead
\toprule
No. & 主張 & 実験と測定されたもの & 出典 & 状態 \\
\midrule\endhead
1 & ヒトの脳のニューロンは約$8.6\cdot10^{10}$個で、\nt{GLUT}ニューロンのほとんどは小脳の小さなニューロンである（表~\ref{tab:types}） & 等方性分画法：組織を溶かして核の均一な懸濁液にし、ニューロンマーカーNeuNをもつ核を数える。成人男性のニューロンは$(86.1\pm8.1)\cdot10^9$個。そのうち大脳皮質にあるのは$19\,\%$にすぎず、大部分は小脳にある & \cite{Azevedo2009} & 測定済み \\
2 & ほとんどすべてのニューロンは主要な神経伝達物質を1種類もつ（命題~\ref{prop:dale}）が、例外もある & マウス脳のトランスクリプトーム・アトラス：約$4\cdot10^6$個の細胞、$5322$クラスタ。合成と小胞充填の遺伝子から各クラスタに神経伝達物質を割り当てた。例外は直接測定されている。スライス標本の\nt{DOPA}ニューロンは同じ小胞トランスポーターを通じてGABAも放出し、線条体のニューロンを抑制する。\nt{DOPA}軸索の一部では、グルタミン酸とドーパミンが同じ配線の異なる終末から出る（電子顕微鏡と光遺伝学） & \cite{Strata1999,Yao2023,Tritsch2012,Zhang2015} & 測定済み \\
3 & 重み行列の列全体が同じ符号をもつ~\eqref{eq:dalesign} & 命題~\ref{prop:dale}と、グルタミン酸受容体はほぼすべて興奮性でGABA受容体は抑制性であることから章内で導出。「1個のニューロンのすべての受け手が同じ符号の重みを受け取る」ことを一般規則として直接検証した例はない。反例は\nt{DOPA}ニューロンによるGABAの共放出（2行目）と、異なる符号の受容体をもつ受け手（11行目） & 章内の導出 & 理論 \\
4 & 大脳皮質のニューロンの4分の3から5分の4が\nt{GLUT}、5分の1から4分の1が\nt{GABA} & サル大脳皮質の$10$領域で、皮質の全層を貫く幅$50$\,\textmu mの柱状領域をGABAで免疫染色。$8$領域ではGABAニューロンが全ニューロンの約$25\,\%$、一次視覚野では$20\,\%$未満 & \cite{Hendry1987} & 測定済み \\
5 & \nt{GLUT}ニューロンの出力は約$10^4$個 & 剖検脳の立体解析学：若年男性の新皮質にシナプス$1.64\cdot10^{14}$個（電子顕微鏡、ダイセクター法）、ニューロン約$2.3\cdot10^{10}$個。比をとるとニューロン1個あたり$\sim7\cdot10^3$シナプス。平均すれば入力数は出力数に等しい & \cite{Tang2001synapses,Pakkenberg1997} & 測定済み \\
6 & \nt{DOPA}ニューロンの出力は$10^5$--$10^6$個の終末に枝分かれする。これは標的の被覆範囲からの推定であって計数ではない & ラットの個々の\nt{DOPA}ニューロンの膜をウイルスで標識し、軸索を完全に再構成：1個のニューロンが線条体の体積の$0.45$から$5.7\,\%$（平均$2.7\,\%$）を覆う。測定したのは被覆範囲で、終末数ではない。終末数は被覆範囲からオーダーとして導いたもので、終末の直接計数はない。\nt{SERO}と\nt{NORA}の終末数は粗い推定 & \cite{Matsuda2009} & 測定済み \\
7 & ブロードキャスト型ニューロンは少ない：\nt{DOPA} $\sim5\cdot10^5$（二つの群のうち主要なほう、下限）、\nt{SERO} $\sim3\cdot10^5$（二つの部分計数の和）、\nt{HIST} $\sim6\cdot10^4$（表~\ref{tab:types}、図~\ref{fig:typepie}） & ヒトでの立体解析学的計数：黒質の色素含有ニューロン$550\,000$個（腹側被蓋野を除く）。背側縫線核のニューロン$235\,000$個（核の全ニューロン）と、橋被蓋のセロトニンニューロンさらに$125\,000$個。視床下部のヒスタミンニューロン約$64\,000$個。\nt{NORA}、\nt{GLYC}、\nt{ACET}、\nt{ADRE}には直接の計数がなく、表の値は粗いオーダー推定である & \cite{Pakkenberg1991,Baker1990,Baker1991,Airaksinen1991} & 測定済み \\
8 & 隙間のグルタミン酸は約1ミリモーラーまで跳ね上がり、$\sim1\ms$で減衰する & シナプス伝達中に、速く解離する競合的阻害薬がNMDA受容体から追い出される程度から求めた（海馬ニューロンの培養）：ピーク$1.1$\,mM、減衰時定数$1.2\ms$ & \cite{Clements1992} & 測定済み \\
9 & 動作中の皮質では興奮性電流と抑制性電流がほぼ釣り合い、ニューロンは閾値付近にいる & 理論：強い結合をもつネットワークは自ら釣り合いの状態に入る。実験：フェレット前頭前皮質のin vivo記録で、「アップ」状態の間、コンダクタンスが$\sim21$\,nS変化してもシナプス電流の反転電位は数百ミリ秒にわたり約$-37$\,mVに保たれる。興奮と抑制が一緒に増減する & \cite{vanVreeswijk1996,Haider2006} & 測定済み \\
10 & \nt{DOPA}ニューロンは報酬予測誤差~\eqref{eq:rpe}を伝える（図~\ref{fig:rpe}） & サルの\nt{DOPA}ニューロンのスパイク：予告なしのジュースにバースト。学習後は合図にバーストし、予測されたジュースには応答しない。約束されたジュースが与えられないと発火率が落ち込む。定量的な実験では、発火率は現在の報酬と過去の報酬の加重平均との差に比例するが、それは正の誤差に限られる。式~\eqref{eq:rpe}そのものは時間差分アルゴリズムである & \cite{Schultz1997,Bayer2005,SuttonBarto2018} & 測定済み \\
11 & 符号と速さを決めるのは受容体：イオンチャネル型は1ミリ秒未満、代謝型ははるかに遅い（命題~\ref{prop:receptor}） & 海馬ニューロンの膜パッチに短いグルタミン酸パルスを与えると、イオンチャネル型受容体の電流は$0.2$--$0.6\ms$で立ち上がり（$20$から$80\,\%$）、$2$--$3\ms$で減衰する。錐体ニューロンの遅い抑制性電位は、GABAの代謝型受容体の選択的阻害薬で消える。ドーパミン受容体には逆の作用をもつ二つのファミリーがある。線条体では、D1受容体をもつニューロンはシナプスの増強にドーパミンを必要とし、D2受容体をもつニューロンは減弱にドーパミンを必要とする & \cite{Colquhoun1992,DutarNicoll1988,Beaulieu2011,Shen2008} & 測定済み \\
12 & ブロードキャストチャネルは1本の配線ではなく、少数の線の束である（\ref{sec:buses}節） & マウス背側縫線核のセロトニンニューロンをウイルスと遺伝学的手法で標識：扁桃体に出力するニューロンと前頭皮質に出力するニューロンは核内の異なる場所にあり、互いに補い合う領域に枝分かれし、不快な刺激に逆向きに応答する。総説：青斑核（\nt{NORA}）のニューロンの亜群は異なる過程を符号化する & \cite{Ren2018,Poe2020} & 測定済み \\
13 & \nt{NORA}がゲイン$g$を決める & 適応的ゲインの仮説。カテコールアミンが伝達特性の傾きを上げるネットワークモデル。ネットワーク全体で「ノルアドレナリンとともに$g$が上がる」ことの直接測定はない。測定されているのは、サルの瞳孔径が青斑核ニューロンのスパイクに追従すること & \cite{AstonJones2005,ServanSchreiber1990,Joshi2016} & 仮説 \\
14 & \nt{ACET}が「書き込み/読み出し」を切り替え、学習率を決める & 仮説。根拠：ラット嗅皮質のスライスで、アセチルコリン作動薬（$100$\,\textmu M）は内部の「想起」シナプスを$60\,\%$以上弱めるが、入力シナプスは$15\,\%$未満しか弱めない & \cite{Hasselmo2006,HasselmoBower1992} & 仮説 \\
15 & \nt{SERO}が割引率$\gamma$を決める。セロトニンニューロンは一様に毎秒数スパイクで発火し、睡眠のある相ではほぼ沈黙する & 「セロトニン$=\gamma$」仮説。最も強い論拠：マウス背側縫線核のセロトニンニューロンを光遺伝学的に活性化すると、遅延報酬を待つのを途中でやめる回数が減る。発火率と急速眼球運動睡眠中の沈黙は測定されている（記録の総説） & \cite{Doya2002,Miyazaki2014,JacobsAzmitia1992} & 仮説 \\
\bottomrule
\end{longtable}}
